\documentclass[11pt,a4paper]{amsart}

\usepackage[utf8]{inputenc}
\usepackage{amsmath,amsthm,amssymb}
\usepackage{graphicx}

\newtheorem{theorem}{Theorem}
\newtheorem{lemma}{Lemma}
\theoremstyle{definition}
\newtheorem{definition}{Definition}

\title{A Monotone Derivative is Continuous}
\author{}
\date{14 March 2026}

\begin{document}

\maketitle

\begin{abstract}
A classical theorem of Darboux states that derivatives have the Intermediate Value Property.
In this note, we prove that if $f'$ is monotone on an interval, then $f'$ is continuous on that interval.
We present two proofs: one based on Darboux's Theorem, and another based on one-sided limits and the Mean Value Theorem.
\end{abstract}

\section{Introduction}

Let $f:I\to\mathbb{R}$ be differentiable on an open interval $I$.
Although the derivative of a differentiable function need not be continuous, it still satisfies the Intermediate Value Property.
For example, the function
\[
f(x)=
\begin{cases}
x^2\sin(1/x), & x\neq 0,\\
0, & x=0
\end{cases}
\]
is differentiable on $\mathbb{R}$, but its derivative is not continuous at $0$.

In this note, we prove the following stronger fact: if $f'$ is monotone on $I$, then $f'$ must in fact be continuous on $I$.

\section{Darboux's Theorem}

We shall use the following classical result.

\begin{theorem}[Darboux]
Let $f$ be continuous on $[a,b]$ and differentiable on $(a,b)$.
If $k$ is a value between $f'(x_1)$ and $f'(x_2)$ for some $x_1,x_2\in(a,b)$ with $x_1<x_2$,
then there exists $c\in(x_1,x_2)$ such that
\[
f'(c)=k.
\]
\end{theorem}

\begin{proof}
Assume $f'(x_1)<k<f'(x_2)$; the other case is similar.
Define
\[
g(x)=f(x)-kx.
\]
Then $g$ is continuous on $[x_1,x_2]$ and differentiable on $(x_1,x_2)$.
Moreover,
\[
g'(x_1)=f'(x_1)-k<0,
\qquad
g'(x_2)=f'(x_2)-k>0.
\]
Since $g$ is continuous on the compact interval $[x_1,x_2]$, it attains a minimum at some point $c\in[x_1,x_2]$.
Because $g'(x_1)<0$, the minimum cannot occur at $x_1$.
Because $g'(x_2)>0$, the minimum cannot occur at $x_2$.
Hence the minimum occurs at some interior point $c\in(x_1,x_2)$.
By Fermat's theorem,
\[
g'(c)=0,
\]
so
\[
f'(c)=k.
\]

\end{proof}

\section{Main Result}

\begin{theorem}
Let $f:I\to\mathbb{R}$ be differentiable on an open interval $I$.
If $f'$ is monotone on $I$, then $f'$ is continuous on $I$.
\end{theorem}

\begin{proof}[Proof 1: Contradiction via Darboux's Theorem]
It suffices to consider the case where $f'$ is monotone increasing; the decreasing case is analogous.

Fix $c\in I$.
Since $f'$ is monotone increasing, the one-sided limits
\[
A:=\lim_{x\to c^-}f'(x),
\qquad
B:=\lim_{x\to c^+}f'(x)
\]
exist.

Suppose that $f'$ is discontinuous at $c$.
Then at least one of the inequalities
\[
A<f'(c), \qquad f'(c)<B
\]
must hold.
In particular, we have
\[
A<B.
\]

Choose $k\in(A,B)$ such that $k\neq f'(c)$.
Now choose points $x_1,x_2\in I$ with
\[
x_1<c<x_2.
\]
By monotonicity,
\[
f'(x_1)\leq A<k<B\leq f'(x_2).
\]
By Darboux's Theorem, there exists $d\in(x_1,x_2)$ such that
\[
f'(d)=k.
\]

We now derive a contradiction.

If $d<c$, then monotonicity gives
\[
f'(d)\leq A<k,
\]
contrary to $f'(d)=k$.

If $d>c$, then monotonicity gives
\[
f'(d)\geq B>k,
\]
again contrary to $f'(d)=k$.

If $d=c$, then
\[
f'(c)=k,
\]
which contradicts the choice of $k$.

Thus $f'$ must be continuous at $c$.
Since $c\in I$ was arbitrary, $f'$ is continuous on $I$.

\end{proof}

\begin{proof}[Proof 2: One-Sided Limits and the Mean Value Theorem]
Again, assume that $f'$ is monotone increasing.

Fix $c\in I$, and let
\[
A:=\lim_{x\to c^-}f'(x),
\qquad
B:=\lim_{x\to c^+}f'(x).
\]
We show that
\[
A=f'(c)=B.
\]

First let $h>0$ be small enough that $c+h\in I$.
By the Mean Value Theorem, there exists $\xi_h\in(c,c+h)$ such that
\[
\frac{f(c+h)-f(c)}{h}=f'(\xi_h).
\]
Since $c<\xi_h<c+h$ and $f'$ is increasing, we have
\[
B\leq f'(\xi_h)\leq f'(c+h).
\]
Hence
\[
B\leq \frac{f(c+h)-f(c)}{h}\leq f'(c+h).
\]
Now let $h\to 0^+$.
Because $f$ is differentiable at $c$,
\[
\lim_{h\to 0^+}\frac{f(c+h)-f(c)}{h}=f'(c),
\]
and by definition of $B$,
\[
\lim_{h\to 0^+}f'(c+h)=B.
\]
Therefore,
\[
B\leq f'(c)\leq B,
\]
so
\[
f'(c)=B.
\]

Next let $h<0$ be close to $0$.
Again by the Mean Value Theorem, there exists $\xi_h\in(c+h,c)$ such that
\[
\frac{f(c+h)-f(c)}{h}=f'(\xi_h).
\]
Since $c+h<\xi_h<c$ and $f'$ is increasing, we have
\[
f'(c+h)\leq f'(\xi_h)\leq A.
\]
Thus
\[
f'(c+h)\leq \frac{f(c+h)-f(c)}{h}\leq A.
\]
Letting $h\to 0^-$, we obtain
\[
A\leq f'(c)\leq A,
\]
hence
\[
f'(c)=A.
\]

Therefore,
\[
A=f'(c)=B,
\]
so $f'$ is continuous at $c$.
Since $c\in I$ was arbitrary, $f'$ is continuous on $I$.

\end{proof}

\end{document}
